% Chapter 8: Automation Pipelines
\chapter{Automation Pipelines}
\label{ch:automation}

Claude Code is not limited to interactive use.
Its headless mode, fan-out capabilities, and the Claude Agent SDK
enable production automation pipelines that run without human interaction.

\section{Headless Mode and CI Integration}
\label{sec:headless}

\marginnote[Official]{\code{claude -p "prompt"} for non-interactive use. \code{--output-format json|stream-json|text}.\sourceurl{https://code.claude.com/docs/en/common-workflows}}

\official{Use \code{claude -p} for non-interactive automation.}
Key flags:

\begin{description}
  \item[\code{-p "prompt"}] Single prompt, no interaction.
  \item[\code{--output-format json}] Structured output for pipeline consumption.
  \item[\code{--allowedTools}] Restrict to specific tools for safety.
  \item[\code{--model}] Override model selection per invocation.
\end{description}

\subsection{Fan-Out Pattern}

\official{For bulk operations across many files:}

\begin{minted}{bash}
# Generate file list
files=$(find src/ -name "*.py" -type f)

# Test on 3 files first
for f in $(echo "$files" | head -3); do
  claude -p "Add type hints to all functions in $f" \
    --allowedTools Read,Edit
done

# Review results, refine prompt if needed, then scale
for f in $files; do
  claude -p "Add type hints to all functions in $f" \
    --allowedTools Read,Edit \
    --output-format json >> results.jsonl
done
\end{minted}

\margintip{Always test on 3--5 files before scaling. The prompt that works on your first file may fail on edge cases in file 47.}

\subsection{Writer/Reviewer Pattern}

\official{Use separate sessions to avoid confirmation bias:}

\begin{enumerate}
  \item \textbf{Writer session}: Makes changes based on requirements.
  \item \textbf{Reviewer session}: Reviews the diff without knowledge of intent.
    ``Review this diff for bugs, regressions, and missed edge cases.
    You have no context about why these changes were made.''
\end{enumerate}

This pattern catches errors that the writer session is blind to
because of its accumulated context about the ``correct'' approach.

\subsection{Permissions in Headless Mode}

\begin{warnbox}[Permissions in CI/CD]
In headless mode (\code{claude -p}), Claude still respects permission rules.
For trusted CI/CD environments, use \code{--dangerously-skip-permissions}
to bypass interactive prompts. This flag should \textbf{only} be used in
environments you fully control (CI runners, Docker containers).
Never expose this flag in user-facing scripts.
\end{warnbox}

% -------------------------------------------------------------------
\section{Claude Agent SDK for Production}
\label{sec:agent-sdk}

\marginnote[Official]{Same agent loop as Claude Code, programmable in Python and TypeScript.\sourceurl{https://platform.claude.com/docs/en/agent-sdk/overview}}

\official{The Claude Agent SDK provides programmatic access to the same
agent loop that powers Claude Code.} It supports:

\begin{itemize}
  \item Filesystem-based configuration (\code{CLAUDE.md}, skills, commands, memory)
    via \code{setting\_sources=["project"]}
  \item Custom tool definitions and MCP server integration
  \item Streaming responses and structured output
  \item Production deployment with error handling and retry logic
\end{itemize}

Use cases:
\begin{description}
  \item[CI/CD Pipelines] Automated code review, test generation, documentation updates.
  \item[Custom Applications] Domain-specific AI assistants with project context.
  \item[Production Automation] Scheduled tasks, monitoring responses, incident triage.
\end{description}

% -------------------------------------------------------------------
\section{Pipeline Design Patterns}
\label{sec:pipeline-patterns}

\practitioner{Production pipelines follow a common architecture:
source $\rightarrow$ extraction $\rightarrow$ validation $\rightarrow$ generation.}

\subsection{Pipeline Directionality}

\begin{warnbox}[Pipeline Direction Violations]
\textbf{Symptom:} Downstream artifacts overwritten by upstream re-runs.

\textbf{Root cause:} Running an extraction script after manually editing
its output. The extraction regenerates from source, destroying your edits.

\textbf{Prevention:} Document flow direction.
Never re-run upstream steps after downstream edits.
Mark generated files with \code{\# GENERATED --- do not edit manually}.
\end{warnbox}

\subsection{Autonomous QA with Traffic-Light Dashboards}

\practitioner{Replace pass/fail with traffic-light thresholds (GREEN/YELLOW/RED).}
This provides nuance: a metric at 78\% coverage is not ``failing''---it is
YELLOW, meaning ``acceptable but track for improvement.''

\begin{minted}{text}
=== Project Health Dashboard ===
Test Coverage:    82% [GREEN]
Lint Warnings:     3  [YELLOW]
Documentation:   91% [GREEN]
Card Quality:    97% [GREEN]
Overall:         HEALTHY (3 GREEN, 1 YELLOW)
\end{minted}

\practitioner{Run dashboards as skills, not manual scripts.}
A skill wraps the dashboard with interpretation:
not just numbers, but recommendations for what to fix first.

\begin{keyinsight}
Automation with Claude follows the same principle as automation without it:
build the pipeline for the common case, handle exceptions manually.
The difference is that Claude can handle a much wider ``common case''
than traditional scripts, reducing the exception rate dramatically.
\end{keyinsight}
